% GENERATED FILE. Edit canonical lessons, then run npm run generate:books.
% canonical-source-hash: fnv1a64:e59194d2ec266a09
% canonical-lessons: SA-C506-vayati, SA-C506-tantuh, SA-C506-ranjakah, SA-C506-ranjayati, SA-C506-saucikah

\chapter{Thread and Loom}
\label{ch:sa-thread-and-loom}
\hlchaptermodality{506}

\begin{chapteropening}
\textbf{By the end of this chapter:} \emph{I can talk about weaving, dyeing and stitching cloth.}

Complete the last lesson of chapter 506: I can talk about weaving, dyeing and stitching cloth.
\end{chapteropening}

\section[\texorpdfstring{\sk{वयति} (vayati)}{वयति (vayati)}]{\sk{वयति} (vayati) — weaves}

\label{lesson:SA-C506-vayati}



\begin{quote}
Before the new one: say the Sanskrit for a kiln, then the Sanskrit for a shape.
\end{quote}

\subsection*{You'll want to know: \sk{वयति}}
\textbf{\sk{वयति}} — \emph{vayati} — \textquotedblleft{}weaves\textquotedblright{}.

\begin{quote}
\sk{तन्तुवायः} \sk{कार्पासस्य} \sk{वस्त्रं} \sk{वयति।}
\end{quote}

\begin{quote}
\emph{tantuvāyaḥ kārpāsasya vastraṃ vayati.}
\end{quote}

\begin{quote}
The weaver weaves cotton cloth.
\end{quote}

Past (he …): \sk{सः} \sk{वस्त्रम्} \sk{उतवान्}, \emph{saḥ vastram utavān}.

\begin{cousinweb}
\textbf{\sk{वयति}} (\emph{vayati}) comes from the root √ve, \textquotedblleft{}to weave\textquotedblright{}, whose present is \emph{vayati}; its past participle is \emph{uta}, and KRM lists the past \emph{utavān} (MW).
\end{cousinweb}

\begin{grammarlens}[title={What you've built}]
One more word for describing what you lived and felt.
\end{grammarlens}

\subsection*{Your turn}
\begin{itemize}
\raggedright
  \item \emph{Say it:} \emph{vayati}
  \item \emph{Say it:} \emph{vayati}, once more
  \item \emph{Say it:} \emph{tantuvāyaḥ kārpāsasya vastraṃ vayati.} — The weaver weaves cotton cloth.
  \item \emph{Recall:} say the Sanskrit for a kiln, then the Sanskrit for a shape, then say \emph{vayati} again
\end{itemize}

\begin{tcolorbox}[breakable,colback=teal!4,colframe=teal!35!black,title={Before you move on}]
What does \sk{वयति} mean? (\textquotedblleft{}Weaves\textquotedblright{}.) Say it once more.
\end{tcolorbox}
\section[\texorpdfstring{\sk{तन्तुः} (tantuḥ)}{तन्तुः (tantuḥ)}]{\sk{तन्तुः} (tantuḥ) — a strand; a thread of the warp}

\label{lesson:SA-C506-tantuh}



\begin{quote}
Before the new one: say the Sanskrit for a shape, then the Sanskrit for weaves.
\end{quote}

\subsection*{You'll want to know: \sk{तन्तुः}}
\textbf{\sk{तन्तुः}} — \emph{tantuḥ} — \textquotedblleft{}a strand; a thread of the warp\textquotedblright{}.

\begin{quote}
\sk{सा} \sk{सूच्यां} \sk{तन्तुं} \sk{स्थापयति।}
\end{quote}

\begin{quote}
\emph{sā sūcyāṃ tantuṃ sthāpayati.}
\end{quote}

\begin{quote}
She threads the needle.
\end{quote}

A masculine noun.

\begin{cousinweb}
\textbf{\sk{तन्तुः}} (\emph{tantuḥ}) comes from the root √tan, \textquotedblleft{}to stretch, to spin out\textquotedblright{}; MW gives \emph{tantu} as a thread, a cord, the warp of a web, a fibre (MW).
\end{cousinweb}

\begin{grammarlens}[title={What you've built}]
One more word for describing what you lived and felt.
\end{grammarlens}

\subsection*{Your turn}
\begin{itemize}
\raggedright
  \item \emph{Say it:} \emph{tantuḥ}
  \item \emph{Say it:} \emph{tantuḥ}, once more
  \item \emph{Say it:} \emph{sā sūcyāṃ tantuṃ sthāpayati.} — She threads the needle.
  \item \emph{Recall:} say the Sanskrit for a shape, then the Sanskrit for weaves, then say \emph{tantuḥ} again
\end{itemize}

\begin{tcolorbox}[breakable,colback=teal!4,colframe=teal!35!black,title={Before you move on}]
What does \sk{तन्तुः} mean? (\textquotedblleft{}A strand; a thread of the warp\textquotedblright{}.) Say it once more.
\end{tcolorbox}
\section[\texorpdfstring{\sk{रञ्जकः} (rañjakaḥ)}{रञ्जकः (rañjakaḥ)}]{\sk{रञ्जकः} (rañjakaḥ) — a dyer}

\label{lesson:SA-C506-ranjakah}



\begin{quote}
Before the new one: say the Sanskrit for weaves, then the Sanskrit for a strand.
\end{quote}

\subsection*{You'll want to know: \sk{रञ्जकः}}
\textbf{\sk{रञ्जकः}} — \emph{rañjakaḥ} — \textquotedblleft{}a dyer\textquotedblright{}.

\begin{quote}
\sk{रञ्जकः} \sk{वस्त्रं} \sk{नीलं} \sk{करोति।}
\end{quote}

\begin{quote}
\emph{rañjakaḥ vastraṃ nīlaṃ karoti.}
\end{quote}

\begin{quote}
The dyer makes the cloth blue.
\end{quote}

A masculine noun.

\begin{cousinweb}
\textbf{\sk{रञ्जकः}} (\emph{rañjakaḥ}) comes from the root √rañj, \textquotedblleft{}to be coloured, to redden\textquotedblright{}, with the agent ending \emph{-aka}; MW gives \emph{rañjaka} as a colourist, a dyer, a painter (MW).
\end{cousinweb}

\begin{grammarlens}[title={What you've built}]
One more word for describing what you lived and felt.
\end{grammarlens}

\subsection*{Your turn}
\begin{itemize}
\raggedright
  \item \emph{Say it:} \emph{rañjakaḥ}
  \item \emph{Say it:} \emph{rañjakaḥ}, once more
  \item \emph{Say it:} \emph{rañjakaḥ vastraṃ nīlaṃ karoti.} — The dyer makes the cloth blue.
  \item \emph{Recall:} say the Sanskrit for weaves, then the Sanskrit for a strand, then say \emph{rañjakaḥ} again
\end{itemize}

\begin{tcolorbox}[breakable,colback=teal!4,colframe=teal!35!black,title={Before you move on}]
What does \sk{रञ्जकः} mean? (\textquotedblleft{}A dyer\textquotedblright{}.) Say it once more.
\end{tcolorbox}
\section[\texorpdfstring{\sk{रञ्जयति} (rañjayati)}{रञ्जयति (rañjayati)}]{\sk{रञ्जयति} (rañjayati) — dyes}

\label{lesson:SA-C506-ranjayati}



\begin{quote}
Before the new one: say the Sanskrit for a strand, then the Sanskrit for a dyer.
\end{quote}

\subsection*{You'll want to know: \sk{रञ्जयति}}
\textbf{\sk{रञ्जयति}} — \emph{rañjayati} — \textquotedblleft{}dyes\textquotedblright{}.

\begin{quote}
\sk{सः} \sk{शाटिकां} \sk{हरिद्रया} \sk{रञ्जयति।}
\end{quote}

\begin{quote}
\emph{saḥ śāṭikāṃ haridrayā rañjayati.}
\end{quote}

\begin{quote}
He dyes the sari with turmeric.
\end{quote}

Past (he …): \sk{सः} \sk{वस्त्रं} \sk{रञ्जितवान्}, \emph{saḥ vastraṃ rañjitavān}.

\begin{cousinweb}
\textbf{\sk{रञ्जयति}} (\emph{rañjayati}) comes from the root √rañj, \textquotedblleft{}to be coloured\textquotedblright{}; MW gives its causative \emph{rañjayati}, to dye, to colour (MW).
\end{cousinweb}

\begin{grammarlens}[title={What you've built}]
One more word for describing what you lived and felt.
\end{grammarlens}

\subsection*{Your turn}
\begin{itemize}
\raggedright
  \item \emph{Say it:} \emph{rañjayati}
  \item \emph{Say it:} \emph{rañjayati}, once more
  \item \emph{Say it:} \emph{saḥ śāṭikāṃ haridrayā rañjayati.} — He dyes the sari with turmeric.
  \item \emph{Recall:} say the Sanskrit for a strand, then the Sanskrit for a dyer, then say \emph{rañjayati} again
\end{itemize}

\begin{tcolorbox}[breakable,colback=teal!4,colframe=teal!35!black,title={Before you move on}]
What does \sk{रञ्जयति} mean? (\textquotedblleft{}Dyes\textquotedblright{}.) Say it once more.
\end{tcolorbox}
\section[\texorpdfstring{\sk{सौचिकः} (saucikaḥ)}{सौचिकः (saucikaḥ)}]{\sk{सौचिकः} (saucikaḥ) — a tailor}

\label{lesson:SA-C506-saucikah}



\begin{quote}
Before the new one: say the Sanskrit for a dyer, then the Sanskrit for dyes.
\end{quote}

\subsection*{You'll want to know: \sk{सौचिकः}}
\textbf{\sk{सौचिकः}} — \emph{saucikaḥ} — \textquotedblleft{}a tailor\textquotedblright{}.

\begin{quote}
\sk{सौचिकः} \sk{सूच्या} \sk{कञ्चुकं} \sk{सीव्यति।}
\end{quote}

\begin{quote}
\emph{saucikaḥ sūcyā kañcukaṃ sīvyati.}
\end{quote}

\begin{quote}
The tailor sews a shirt with a needle.
\end{quote}

A masculine noun.

\begin{cousinweb}
\textbf{\sk{सौचिकः}} (\emph{saucikaḥ}) is built on \emph{sūci}, \textquotedblleft{}a needle\textquotedblright{}, with the ending \emph{-ika}; MW gives \emph{saucika} as one who lives by his needle, a tailor (MW).
\end{cousinweb}

\begin{grammarlens}[title={What you've built}]
One more word for describing what you lived and felt.
\end{grammarlens}

\subsection*{Your turn}
\begin{itemize}
\raggedright
  \item \emph{Say it:} \emph{saucikaḥ}
  \item \emph{Say it:} \emph{saucikaḥ}, once more
  \item \emph{Say it:} \emph{saucikaḥ sūcyā kañcukaṃ sīvyati.} — The tailor sews a shirt with a needle.
  \item \emph{Recall:} say the Sanskrit for a dyer, then the Sanskrit for dyes, then say \emph{saucikaḥ} again
\end{itemize}

\begin{tcolorbox}[breakable,colback=teal!4,colframe=teal!35!black,title={Before you move on}]
What does \sk{सौचिकः} mean? (\textquotedblleft{}A tailor\textquotedblright{}.) Say it once more.
\end{tcolorbox}
